\section{Why the known-target oracle cannot be the general DC-SBM benchmark}
\label{sec:impossibility}

The profile reference and an oracle supplied with the target's true degree
parameter are different experiments. The distinction cannot be removed by
a more accurate spectral reconstruction: the following obstruction holds
even when the classifier observes the entire adjacency matrix. The example
has bounded positive degree parameters, exact block normalization, balanced
groups, a positive definite block matrix, and logarithmic expected degree.
Every probability calculation below is for the Bernoulli graph itself.

\paragraph{The local known-target reference.}
For a fixed DC-SBM parameter point \(\tau=(z^\tau,\theta^\tau)\), write
\(\PP_\tau\) for the graph law. At a target vertex \(v\), give a local oracle
the values of \(p,C,\theta_v^\tau\) and all neighbor parameters
\((z_j^\tau,\theta_j^\tau)\), \(j\ne v\). Its candidate row laws are
\begin{equation}
 q_{vj}^{(a,\tau)}
   =p\theta_v^\tau\theta_j^\tau C_{a z_j^\tau},\qquad
 L_v^{(a,\tau)}(A_{v,-v})
   =\prod_{j\ne v}
     \bigl(q_{vj}^{(a,\tau)}\bigr)^{A_{vj}}
     \bigl(1-q_{vj}^{(a,\tau)}\bigr)^{1-A_{vj}}.
 \label{eq:imposs-local-laws}
\end{equation}
Let \(\delta_v^{\mathrm{known},\tau}\) maximize these two likelihoods with
equal candidate priors and a fixed tie rule, and define
\begin{equation}
 R_{\mathrm{known}}^\tau(n)
   =\frac1n\sum_{v=1}^n
     \PP_\tau\{\delta_v^{\mathrm{known},\tau}\ne z_v^\tau\}.
 \label{eq:imposs-known-risk}
\end{equation}
The target parameter is held at its supplied true value under both candidate
labels. The local comparison does not impose a global exact-group-size or
block-normalization constraint after revealing the other labels; an exact
group-size constraint would itself determine a missing label. Thus
\eqref{eq:imposs-known-risk} is the nondegenerate local likelihood benchmark,
not a globally constrained oracle experiment.

\begin{theorem}[An exact Bernoulli obstruction for the known-target oracle]
\label{thm:imposs}
Let \(n=2m\) tend to infinity through even integers, and set
\begin{equation}
 C=\begin{pmatrix}4&2\\2&11/10\end{pmatrix},\qquad
 p=\frac{M\log n}{n},\qquad M=100.
 \label{eq:imposs-model}
\end{equation}
There are two explicit parameter sequences, denoted \(+\) and \(-\), for
the independent-edge Bernoulli DC-SBM with probabilities
\(p\theta_u\theta_vC_{z_uz_v}\), such that both sequences satisfy
\begin{equation}
 n_1=n_2=m,\qquad
 \sum_{v:z_v=a}\theta_v=m\quad(a=1,2),\qquad
 \frac12\le\theta_v\le\frac32.
 \label{eq:imposs-normalization}
\end{equation}
For every possibly randomized clustering algorithm based on the full graph
and information common to the two parameter points, put
\begin{equation}
 R_{\mathrm{alg}}^\tau(n)
   =\E_\tau\left[
       \frac1n\min_{\sigma\in\mathfrak S_2}
       \sum_{v=1}^n\ind{\sigma(\widehat z_v)\ne z_v^\tau}
     \right],\qquad \tau\in\{+,-\}.
 \label{eq:imposs-alg-risk}
\end{equation}
Define
\begin{equation}
 c=\frac{(2-\sqrt2)^2}{8},\qquad
 \gamma=\frac{(\sqrt{11/10}-1)^2}{2}.
 \label{eq:imposs-constants}
\end{equation}
Then
\begin{equation}
 \max_{\tau\in\{+,-\}}
 \frac{R_{\mathrm{alg}}^\tau(n)}{R_{\mathrm{known}}^\tau(n)}
 \ge \frac{1+o(1)}4\,
       n^{M(c-2\gamma)-1}.
 \label{eq:imposs-ratio}
\end{equation}
The exponent is strictly positive; at \(M=100\),
\begin{equation}
 M(c-2\gamma)-1
   =3.05109151537556\ldots>0.
 \label{eq:imposs-positive-gap}
\end{equation}
In particular, no algorithm with unknown vertex degree parameters can have
a bounded risk ratio to \eqref{eq:imposs-known-risk} along every admissible
parameter sequence in this class. This includes every algorithm based only
on retained spectral information and observed degrees.
\end{theorem}

\begin{proof}
\textbf{Construction and exact normalization.}
Distinguish vertices \(i=1\) and \(k=2\). Partition the remaining vertices
into two fixed anchor sets \(V_1,V_2\), each of size \(m-1\). Every anchor
in \(V_a\) has label \(a\) under both parameter points. Give anchors in
\(V_1\) the common degree parameter
\[
 t_m=1+\frac1{2(m-1)},
\]
and anchors in \(V_2\) degree parameter one. At the two parameter points set
\begin{equation}
 \begin{aligned}
 +:\quad &(z_i,\theta_i)=(1,1/2),
          & (z_k,\theta_k)&=(2,1),\\
 -:\quad &(z_i,\theta_i)=(2,1),
          & (z_k,\theta_k)&=(1,1/2).
 \end{aligned}
 \label{eq:imposs-two-points}
\end{equation}
Each group has exactly \(m\) vertices. Its total degree parameter is exactly
\(m\), since
\[
 (m-1)t_m+\frac12=m,
 \qquad (m-1)+1=m.
\]
The stated degree bounds hold for \(m\ge2\). The matrix \(C\) has positive
entries, positive trace, and determinant \(2/5>0\), so it is positive
definite with a fixed positive smallest eigenvalue. All edge probabilities
belong to \((0,1)\) for sufficiently large \(n\).

\medskip
\textbf{Only \(n-2\) Bernoulli coordinates distinguish the graph laws.}
All anchor--anchor probabilities agree between the two parameter points.
The edge \(i\)--\(k\) has probability \(p\) under either point. For an edge
from either distinguished vertex to an anchor of \(V_1\), the two candidate
intensities agree exactly:
\[
 \frac12\cdot4\cdot t_m=1\cdot2\cdot t_m=2t_m.
\]
Thus the only distinct coordinates are the \(2(m-1)=n-2\) edges from
\(\{i,k\}\) to \(V_2\). Their two probabilities are \(p\) and
\((11/10)p\), exchanged between \(+\) and \(-\).

For probability mass functions \(P,Q\), let
\(\mathcal A(P,Q)=\sum_x\sqrt{P(x)Q(x)}\) be their affinity. Independence
gives the exact graph affinity
\begin{equation}
 \mathcal A_n
 =\mathcal A(\PP_+,\PP_-)
 =\left[
     \sqrt{11/10}\,p+
       \sqrt{(1-p)(1-(11/10)p)}
   \right]^{n-2}.
 \label{eq:imposs-affinity}
\end{equation}
The factor in brackets has the elementary expansion
\[
 \sqrt{11/10}\,p+
       \sqrt{(1-p)(1-(11/10)p)}
 =1-\gamma p+O(p^2).
\]
Since \(np^2=M^2(\log n)^2/n=o(1)\),
\begin{equation}
 \mathcal A_n^2
   =\exp\{-2M\gamma\log n+o(1)\}
   =(1+o(1))n^{-2M\gamma}.
 \label{eq:imposs-affinity-rate}
\end{equation}
The elementary Cauchy--Schwarz bound
\(\TV(P,Q)\le\sqrt{1-\mathcal A(P,Q)^2}\) implies that every test \(\phi\)
of the two graph laws satisfies
\begin{align}
 \frac{\PP_+(\phi=-)+\PP_-(\phi=+)}2
 &\ge \frac{1-\TV(\PP_+,\PP_-)}2 \\
 &\ge \frac{\mathcal A_n^2}{4}
  =\frac{1+o(1)}4 n^{-2M\gamma},
 \label{eq:imposs-test-lower}
\end{align}
where the second inequality uses
\(1-\sqrt{1-y}\ge y/2\). Independent randomization by the test or the
clustering algorithm does not alter this bound: adjoining a common random
seed leaves the affinity unchanged.

\medskip
\textbf{From testing to graph-average permutation loss.}
Orient the two output labels of an arbitrary estimated partition by the
permutation minimizing disagreement with the fixed anchor labels, using
a fixed tie rule. This is one measurable rule under both parameter points,
since the anchors and their labels are the same. Define a test \(\phi\) to
declare \(+\) if the oriented label of vertex \(i\) is 1, and \(-\)
otherwise.

If the estimated partition has zero permutation-invariant error, the
anchor orientation recovers its correct permutation uniquely: both anchor
sets are nonempty. Therefore an erroneous test outcome implies at least
one error in the permutation-invariant count. Pathwise, under either point,
\begin{equation}
 \frac1n\min_{\sigma\in\mathfrak S_2}
       \sum_v\ind{\sigma(\widehat z_v)\ne z_v^\tau}
 \ge\frac1n\ind{\phi\ne\tau}.
 \label{eq:imposs-orientation}
\end{equation}
Taking expectations and using \eqref{eq:imposs-test-lower} gives
\begin{equation}
 \frac{R_{\mathrm{alg}}^+(n)+R_{\mathrm{alg}}^-(n)}2
 \ge\frac{1+o(1)}4 n^{-2M\gamma-1}.
 \label{eq:imposs-graph-risk-lower}
\end{equation}
The anchors are needed only to turn a partition into a test in this proof.
Indeed, their labels and degree parameters, the block matrix, the exact
group sizes, and the block-normalization constraints could all be revealed
to the clustering algorithm: they are common to the two parameter points,
and \eqref{eq:imposs-test-lower} would still apply.

\medskip
\textbf{An upper bound on every known-target local error.}
For any target \(v\), at least \(m-2\) anchors in \(V_1\) are different
from \(v\). For each such anchor \(j\), its two local candidate edge
probabilities in \eqref{eq:imposs-local-laws} are
\[
 u=4p\theta_v\theta_j,
 \qquad v'=2p\theta_v\theta_j.
\]
They satisfy
\[
 (\sqrt u-\sqrt{v'})^2
 =p\theta_v\theta_j(2-\sqrt2)^2
 \ge\frac p2(2-\sqrt2)^2.
\]
For any Bernoulli probabilities \(u,v'\), the arithmetic--geometric mean
inequality for their nonedge probabilities yields
\begin{equation}
 \sqrt{uv'}+\sqrt{(1-u)(1-v')}
 \le 1-\frac{(\sqrt u-\sqrt{v'})^2}{2}
 \le \exp\left\{-\frac{(\sqrt u-\sqrt{v'})^2}{2}\right\}.
 \label{eq:imposs-one-edge-upper}
\end{equation}
Under either local true label, Markov's inequality applied to the square
root of the competing likelihood ratio bounds the LRT's directional error
by the product affinity of all its row coordinates. Discarding coordinates
outside these \(m-2\) anchors can only increase that upper bound. Hence,
uniformly over targets and both parameter points,
\begin{align}
 \PP_\tau\{\delta_v^{\mathrm{known},\tau}\ne z_v^\tau\}
 &\le\exp\left\{-\frac{(m-2)p}{4}(2-\sqrt2)^2\right\} \\
 &= (1+o(1))n^{-Mc}.
 \label{eq:imposs-oracle-upper}
\end{align}
In particular,
\begin{equation}
 R_{\mathrm{known}}^+(n)+R_{\mathrm{known}}^-(n)
 \le 2(1+o(1))n^{-Mc}.
 \label{eq:imposs-oracle-sum-upper}
\end{equation}

Finally, the maximum of two nonnegative risk ratios is at least their
denominator-weighted average. Combining
\eqref{eq:imposs-graph-risk-lower} and
\eqref{eq:imposs-oracle-sum-upper} therefore proves
\eqref{eq:imposs-ratio}. The constants in
\eqref{eq:imposs-constants} are
\[
 c=0.04289321881345246\ldots,
 \qquad \gamma=0.00119115182984846\ldots,
\]
which give the exponent in \eqref{eq:imposs-positive-gap}. Its positivity
does not depend on numerical evaluation: \(c>1/25\) and
\(2\gamma<1/400\), so \(100(c-2\gamma)-1>11/4>0\).
\end{proof}

\begin{remark}[Scope of the obstruction]
\label{rem:imposs-scope}
The same proof works for every fixed
\(M>(c-2\gamma)^{-1}\), so its logarithmic-degree multiplier can exceed
any prescribed fixed lower threshold. The example also satisfies
\(np^2\to0\). Thus neither a larger logarithmic expected degree nor the
sparse Bernoulli approximation removes the nuisance-information gap.

The result concerns a uniform comparison to the known-target reference,
not the full-edge profile benchmark used by the positive DC-SBM theorem.
It also does not assert that every fixed DC-SBM has an unbounded ratio:
it proves that at least one of the two explicit parameter sequences has
an unbounded limsup ratio for any algorithm. Additional restrictions can
make the stronger benchmark attainable. For example, the balanced block
strengths in this counterexample are
\(C(1/2,1/2)^T=(3,31/20)^T\), which are unequal; the example lies outside
the equal-exposure corollary.

The counterexample is also outside the assortative parameter class of
Gao, Ma, Zhang and Zhou~\cite{gmzz2018}: its smallest diagonal entry is
$1.1$, whereas its off-diagonal entry is $2$. Their class requires every
within-group strength to exceed every between-group strength. Therefore
this obstruction is not a negative result for their classical balanced
planted-partition benchmark. That benchmark is explicitly covered by
\cref{prop:classical-dc-oracle-bridge} on the stated bounded-parameter
class. Their classwise minimax theorem and a pointwise comparison to a
known-target oracle on every arbitrary block matrix are different claims.

Finally, a lower bound for an experiment retaining only block counts would
not by itself establish a lower bound for arbitrary data-dependent real
spectral coordinates. The present theorem needs no such compression claim:
it is a lower bound for the full adjacency experiment with unknown target
degree parameters.
\end{remark}
